\section{Metodología}

\subsection{Diseño de investigación}

DSR busca construir y evaluar artefactos relevantes con rigor \cite{hevner2004dsr}. El estudio sigue las seis actividades de Peffers \emph{et al.}: identificar la necesidad, definir objetivos, diseñar y desarrollar, demostrar, evaluar y comunicar \cite{peffers2007dsrm}. La Fig.~\ref{fig:dsr} muestra cómo cada evaluación alimenta una nueva iteración. Los cuadernos preliminares de Qwen con cinco categorías de moderación y de Transformers jerárquicos que no producen resultados se registran como diseños abandonados, no como experimentos ejecutados.

\begin{figure*}[!t]
\centering
\resizebox{\textwidth}{!}{\begin{tikzpicture}[
  font=\scriptsize,
  phase/.style={draw,rounded corners=2pt,minimum width=21mm,minimum height=11mm,align=center,fill=blue!7},
  evidence/.style={draw,rounded corners=2pt,text width=20mm,minimum height=12mm,align=center,fill=green!9},
  arr/.style={-{Stealth[length=1.8mm]},thick}
]
\node[phase] (p1) {1. Necesidad\\y requisitos};
\node[phase,right=3mm of p1] (p2) {2. Objetivos\\medibles};
\node[phase,right=3mm of p2] (p3) {3. Diseño y\\construcción};
\node[phase,right=3mm of p3] (p4) {4. Demostración\\en videos};
\node[phase,right=3mm of p4] (p5) {5. Evaluación\\sin fuga};
\node[phase,right=3mm of p5] (p6) {6. Comunicación\\y despliegue};
\foreach \a/\b in {p1/p2,p2/p3,p3/p4,p4/p5,p5/p6}{\draw[arr] (\a)--(\b);}
\node[evidence,below=6mm of p1] (e1) {Corpus\\versionado};
\node[evidence,right=3mm of e1] (e2) {Clásicos de\\cinco categorías};
\node[evidence,right=3mm of e2] (e3) {Consolidación y\\adjudicación};
\node[evidence,right=3mm of e3] (e4) {MiniLM/E5\\cinco categorías};
\node[evidence,right=3mm of e4] (e5) {14 configuraciones\\comparadas};
\node[evidence,right=3mm of e5] (e6) {Frontend 05\\y retroalimentación};
\foreach \a/\b in {e1/e2,e2/e3,e3/e4,e4/e5,e5/e6}{\draw[arr] (\a)--(\b);}
\draw[arr,dashed] (p5.south) |- ++(0,-2mm) -| (p3.south);
\node[below=1mm of e3,align=center,text width=.95\textwidth] {\textit{Cada evaluación devolvió evidencia a la siguiente iteración; un resultado no ejecutado se conserva como plan, no como hallazgo.}};
\end{tikzpicture}
}
\caption{Ciclo DSR y evidencia producida por las iteraciones del proyecto. Elaboración propia basada en las seis actividades de Peffers \emph{et al.} \cite{peffers2007dsrm}.}
\label{fig:dsr}
\end{figure*}

\subsection{Fuentes y criterios de selección}

El universo de adquisición comprende videos públicos de YouTube vinculados con el contexto peruano. La búsqueda abarca registros políticos, periodísticos, humorísticos, de farándula, entretenimiento, deportes, gastronomía y viajes para no confundir un único estilo con daño. Un video es elegible si tiene subtítulos en español, al menos 200 caracteres útiles y metadatos suficientes para identificarlo. El corpus utilizado por los experimentos se reporta siempre mediante su total integrado final.

El corpus construido y utilizado en los experimentos corresponde al esquema 2.0 y al corte del 27 de julio de 2026. La Tabla~\ref{tab:fuentes} presenta sus totales finales. El Apéndice~\ref{app:fuentes_dataset} detalla las fuentes, la procedencia de las etiquetas y los canales. Un manifiesto registra el hash SHA-256 del corpus, según el estándar de hash seguro \cite{nist2015sha}, y la información necesaria para reconstruir la procedencia de cada fila. El corpus todavía no tiene DOI ni una licencia que autorice la redistribución de los textos.

La adquisición combina un catálogo de canales con búsquedas dirigidas a clases minoritarias. La muestra es intencional: prioriza fuentes con mayor rendimiento previo de categorías de moderación y consultas temáticas. La selección informativa y el balance en aprendizaje activo son antecedentes generales \cite{settles2009active,fairstein2024balancing}; la regla concreta de búsqueda es una heurística del proyecto. El proceso excluye videos ya procesados, videos sin subtítulos, textos demasiado cortos, fallas de extracción y duplicados por \texttt{video\_id} o hash normalizado. El análisis utiliza texto y metadatos, sin descargar audio ni imagen. La Tabla~\ref{tab:fuentes} reporta el corpus integrado final como una sola unidad.

\begin{table}[t]
\centering
\caption{Tamaño y distribución del corpus integrado final}
\label{tab:fuentes}
\small
\begin{tabular}{@{}lr@{}}
\toprule
\textbf{Componente} & \textbf{Total} \\
\midrule
Videos & 3\,230 \\
\multicolumn{2}{@{}c@{}}{\makebox[69mm]{\leaders\hbox{\rule{4pt}{0.35pt}\hspace{2.5pt}}\hfill\kern0pt}} \\[-0.4ex]
Chunks (máx. observado: 170 palabras) & 117\,244 \\
\multicolumn{2}{@{}c@{}}{\makebox[69mm]{\leaders\hbox{\rule{4pt}{0.35pt}\hspace{2.5pt}}\hfill\kern0pt}} \\[-0.4ex]
\multicolumn{2}{@{}l}{\textbf{Etiquetas}} \\
\hspace{1em}Seguro & 110\,181 \\
\hspace{1em}Algún daño & 7\,063 \\
\hspace{2em}\Rac & 1\,856 \\
\hspace{2em}\Gen & 1\,961 \\
\hspace{2em}\Ame & 3\,028 \\
\hspace{2em}\Sex & 2\,282 \\
\bottomrule
\end{tabular}
\begin{minipage}{.98\columnwidth}\vspace{1mm}\footnotesize
\emph{Nota:} el máximo observado cuenta palabras separadas por espacios en el corpus final; la segmentación se cierra a los 30 segundos o 600 caracteres y no fija un máximo de palabras. Las cuatro categorías subordinadas son incidencias multietiqueta y no suman “algún daño”. \Ame{} es la unión de acoso personal y amenaza directa; 514 chunks contienen ambos subtipos. Los conteos proceden del corpus canónico.
\end{minipage}
\end{table}

El Apéndice~\ref{app:fuentes_dataset} desagrega las categorías de fuente y la procedencia de las etiquetas. También presenta individualmente los canales con más de dos videos y agrupa los canales con uno o dos, cuyas direcciones permanecen en los metadatos del corpus.

\subsection{Subtítulos, limpieza y chunks}

La fuente textual comprende pistas de subtítulos públicas, manuales o automáticas según disponibilidad. La extracción usa VTT con \texttt{yt-dlp} \cite{ytdlp2026} y, cuando hace falta, YouTube Transcript API \cite{depoix2026transcript}. El sistema intenta la familia \texttt{es-PE}/\texttt{es-419}/\texttt{es}; los manifiestos preservan el método, los identificadores y los hashes del resultado. El proyecto usa las pistas existentes y no realiza reconocimiento automático de voz propio.

La limpieza aplica normalización Unicode NFKC \cite{unicode2025normalization}, unifica espacios y ajusta texto y tiempos. Antes de agrupar, elimina hasta 12 palabras repetidas entre segmentos VTT consecutivos. Después cierra cada chunk al alcanzar 30 segundos o 600 caracteres y exige al menos 90 caracteres. La regla de segmentación es propia del proyecto y genera chunks sin solapamiento intencional. Tras deduplicar y retirar decisiones no elegibles, el único total reportado para el corpus es el de la Tabla~\ref{tab:fuentes}.

\subsection{Taxonomía y evolución del contrato}

El archivo \path{archivo/taxonomia_v1_3/para_equiquetado_LLM} conserva el paquete histórico de anotación manual o asistida: guía v1.3, taxonomía, instrucciones, entrada y ejemplos de salida. El vocabulario fino y los criterios operativos coinciden con las copias auditadas del proyecto. Para los lotes anotados por API, el sistema combinó el prompt compacto v1.1, la taxonomía completa y las reglas técnicas de salida; sus manifiestos conservan la versión del prompt. El Anexo~\ref{app:prompt_operativo} reproduce ese texto y delimita las corridas que acredita. La auditoría muestra 14 etiquetas finas: 12 fenómenos de daño y dos estados seguros, además de tres flags transversales. La separación entre blanco individual/grupal y daño explícito/implícito sigue el marco de Waseem \emph{et al.} \cite{waseem2017abuse}; las definiciones generales de ataque por identidad, doxeo, insulto, agresión sexual y amenaza se contrastan con Banko \emph{et al.} \cite{banko2020taxonomy}. La Tabla~\ref{tab:taxonomia} documenta la adaptación y sus límites. No se atribuye a esos autores la lista exacta del proyecto.

La Fig.~\ref{fig:datos} sintetiza la relación entre categorías operativas, fenómenos finos, estados seguros y señales de revisión. La Fig.~\ref{fig:pipeline} enlaza esa taxonomía con la adquisición, la preanotación, la adjudicación y la evaluación; ambas se sitúan aquí para acompañar la primera explicación del método.

\begin{figure*}[!t]
\centering
\resizebox{.96\textwidth}{!}{\begin{tikzpicture}[
  font=\scriptsize,
  coarse/.style={draw,rounded corners=2pt,fill=blue!9,text width=34mm,minimum height=11mm,align=center,inner sep=3pt},
  fine/.style={draw,rounded corners=2pt,fill=orange!9,text width=34mm,minimum height=27mm,align=left,inner sep=3pt},
  safe/.style={draw,rounded corners=2pt,fill=green!10,align=center,minimum height=12mm,inner sep=3pt},
  flag/.style={draw,rounded corners=2pt,fill=purple!9,align=center,minimum height=15mm,inner sep=3pt},
  root/.style={draw,rounded corners=2pt,fill=gray!10,text width=47mm,minimum height=9mm,align=center},
  arr/.style={-{Stealth[length=1.8mm]},thick},
  bus/.style={thick}
]
% Capa 1: contrato activo. El bus queda en un carril vacío por encima de las cajas.
\node[root] (contract) at (6,2.05) {\textbf{Contrato activo multietiqueta}\\cuatro salidas de daño};
\node[coarse] (rac) at (0,0) {\textbf{RACISMO /\\DISCRIMINACIÓN}};
\node[coarse] (gen) at (4,0) {\textbf{GÉNERO}};
\node[coarse] (ame) at (8,0) {\textbf{ACOSO / AMENAZA}\\OR de dos salidas preliminares};
\node[coarse] (sex) at (12,0) {\textbf{CONTENIDO\\SEXUAL}};
\coordinate (topbus) at (6,1.10);
\draw[bus] (contract.south) -- (topbus);
\draw[bus] (0,1.10) -- (12,1.10);
\foreach \n/\x in {rac/0,gen/4,ame/8,sex/12}{
  \draw[arr] (\x,1.10) -- (\n.north);
}

% Capa 2: fenómenos finos, contenidos en la columna de su salida gruesa.
\node[fine] (rf) at (0,-2.35) {\textbf{5 fenómenos de daño}\\
  étnico explícito\\lingüístico\\clasismo racial\\regional\\encubierto};
\node[fine] (gf) at (4,-2.35) {\textbf{2 fenómenos de daño}\\
  misoginia / acoso de género\\homofobia / transfobia};
\node[fine] (af) at (8,-2.35) {\textbf{2 fenómenos de daño}\\
  ataque dirigido / acoso personal\\amenaza directa};
\node[fine] (sf) at (12,-2.35) {\textbf{3 fenómenos de daño}\\
  sexual explícito\\cosificación sexual\\material no consensual};
\foreach \a/\b in {rac/rf,gen/gf,ame/af,sex/sf}{
  \draw[arr] (\a.south) -- (\b.north);
}

% Capa 3: estados seguros y señales transversales en carriles separados.
\node[safe,text width=46mm] (safestates) at (.60,-4.85) {\textbf{2 estados seguros}\\
  \texttt{seguro}; \texttt{seguro\_ironia\_marcada}\\
  mutuamente excluyentes con daño};
\node[safe,text width=24mm] (safeout) at (4.50,-4.85) {\textbf{SEGURO derivado}\\si ninguna de 4\\salidas está activa};
\draw[arr] (safestates.east) -- (safeout.west);

\node[flag,text width=55mm] (flags) at (8.95,-4.85) {\textbf{3 flags transversales}\\
  \texttt{ironia\_ambigua}; \texttt{humor\_encubridor};\\
  \texttt{contexto\_necesario}\\
  acompañan un daño; no son categorías de moderación};
\node[flag,text width=28mm,fill=yellow!14] (human) at (13.70,-4.85) {\textbf{Revisión humana}\\obligatoria según\\el protocolo};
\draw[arr,dashed] (flags.east) -- (human.west);
\end{tikzpicture}
}
\caption{Contrato activo v2: \texttt{SEGURO} y cuatro categorías de daño se entrenan como cinco salidas; las 14 etiquetas finas y tres señales transversales se conservan. La fusión de acoso personal y amenaza es una decisión operativa, no una equivalencia legal. Elaboración propia informada por tipologías de abuso y contenido dañino \cite{waseem2017abuse,banko2020taxonomy} y por fuentes peruanas citadas en la Tabla~\ref{tab:taxonomia}.}
\label{fig:datos}
\end{figure*}

\begin{figure*}[!t]
\centering
\resizebox{\textwidth}{!}{\begin{tikzpicture}[
  x=1cm,y=1cm,font=\scriptsize,
  box/.style={draw,rounded corners=2pt,text width=27mm,minimum height=12mm,align=center,inner sep=3pt,fill=blue!7},
  model/.style={box,fill=orange!10},
  human/.style={box,fill=green!10},
  arr/.style={-{Stealth[length=1.8mm]},thick},
  futurearr/.style={arr,dashed},
  bus/.style={thick}
]
\node[box] (source) at (0,0) {Videos públicos\\con subtítulos};
\node[box] (download) at (3.15,0) {\texttt{yt-dlp}\\subtítulo manual\\o automático};
\node[box] (chunk) at (6.30,0) {Limpieza y chunks\\30 s / 600 car.\\depura repetición VTT};
\node[model] (flash) at (9.45,0) {Preanotación\\DeepSeek Flash};
\node[model] (pro) at (12.60,0) {Revisión dirigida\\DeepSeek Pro};
\node[human] (human) at (15.75,0) {Adjudicación final\\asistida};
\foreach \a/\b in {source/download,download/chunk,chunk/flash,flash/pro,pro/human}{
  \draw[arr] (\a)--(\b);
}

\node[box] (dataset) at (9.45,-2.30) {Corpus versionado\\117\,244 chunks\\partición por video};
\node[model] (train) at (12.60,-2.30) {Entrenamiento\\clásicos, MiniLM/E5\\y Qwen3--LoRA};
\node[human] (audit) at (15.75,-2.30) {Validación y test\\separados;\\auditoría fina};
\coordinate (databus) at (12.60,-1.22);
\draw[bus] (flash.south) -- (flash.south |- databus);
\draw[bus] (flash.south |- databus) -- (human.south |- databus);
\draw[bus] (pro.south) -- (pro.south |- databus);
\draw[bus] (human.south) -- (human.south |- databus);
\draw[arr] (flash.south |- databus) -- (dataset.north);
\draw[arr] (dataset)--(train);
\draw[arr] (train)--(audit);

\node[human,text width=39mm] (future) at (15.75,-4.35) {Revisión operativa futura\\nueva cohorte versionada};
\draw[futurearr] (audit.south)--(future.north);
\end{tikzpicture}
}
\caption{Flujo de adquisición, preanotación, adjudicación y evaluación. La evaluación no retorna al corpus evaluado; la revisión operativa solo puede originar una versión futura y separada. Elaboración propia.}
\label{fig:pipeline}
\end{figure*}

El contrato usado para los resultados de este artículo contiene cuatro salidas de daño: \Rac, \Gen, \Ame{} y \Sex. \Ame{} une acoso personal y amenaza directa mediante OR/máximo: la amenaza directa tiene menor soporte (688 incidencias frente a 2\,854 de acoso personal), de modo que la unión aumenta el soporte, pero no establece equivalencia semántica o legal. En ese baseline, \texttt{SEGURO} se derivó cuando ninguna salida de daño cruzó su umbral.

El contrato reproducible v2 corrige esa asimetría: entrena cinco salidas desde el inicio, \texttt{SEGURO}, \Rac, \Gen, \Ame{} y \Sex. \texttt{SEGURO} tiene ejemplos, score, umbral y métricas propios y es mutuamente excluyente con cualquier daño; las cuatro salidas de daño siguen siendo multietiqueta. Un conflicto \texttt{SEGURO}+daño, ninguna salida o un score próximo al umbral obliga revisión. Los casos indeterminados usan \texttt{needs\_review=true} y \texttt{training\_eligible=false}; no forman una sexta clase ni se convierten en seguros. Esta revisión de contrato no altera retrospectivamente las métricas del baseline: requiere una cabeza nueva, reentrenamiento y recalibración.

\begin{table*}[t]
\centering
\caption{Definición y sustento de la taxonomía operativa multietiqueta}
\label{tab:taxonomia}
\small
\begin{tabularx}{\textwidth}{>{\raggedright\arraybackslash}p{27mm}>{\raggedright\arraybackslash}p{45mm}>{\raggedright\arraybackslash}X>{\raggedright\arraybackslash}p{38mm}}
\toprule
\textbf{Salida / señal} & \textbf{Etiquetas finas} & \textbf{Inclusión operacional y exclusiones} & \textbf{Antecedentes que orientan la decisión} \\
\midrule
\Rac & \texttt{racismo\_etnico\_explicito}, \texttt{racismo\_linguistico}, \texttt{clasismo\_racial}, \texttt{discriminacion\_regional}, \texttt{racismo\_encubierto} & Ataque o inferiorización por identidad étnico-racial, origen regional, variedad lingüística o clase racializada; incluye formas encubiertas tras criterios de educación/cultura. Una mención identitaria, el español andino o una crítica no degradante no bastan. & Callirgos; Zavala y Zariquiey; Zavala y Back; Brañez; Zavala y Almeida; Salem; Portocarrero/Vich \cite{callirgos1993racismo,zavala2007discurso,zavala2017racismo,branez2012amixer,almeida2022motoso,salem2016amixer,portocarrero2009racismo,vich2018dinamicas}. \\
\Gen & \texttt{misoginia\_acoso\_genero}, \texttt{homofobia\_transfobia} & Ataque, degradación sexualizada o exclusión por ser mujer, por género, orientación sexual o identidad de género. No activa por informar, apoyar o mencionar a una población. & EXIST; Zeinert \emph{et al.}; informes peruanos; Chakravarthi; Rottenbacher; Lovón-Cueva y Lovón-Cueva \cite{rodriguez2021exist,zeinert2021misogyny,albornoz2018conocer,defensoria2021violenciaenlinea,chakravarthi2024homotrans,rottenbacher2012homofobia,lovon2022lesbofobia}. \\
\Ame & \texttt{acoso\_personal}, \texttt{amenaza\_directa} & Ataque dirigido a una persona identificable o exposición de datos privados; expresión de intención de daño físico y, por regla local, legal o económico. Una crítica aislada no degradante o una noticia que cita una amenaza sin respaldarla requiere contexto y no activa automáticamente. & Ataque dirigido, doxeo y amenaza \cite{waseem2017abuse,wulczyn2017exmachina,banko2020taxonomy,defensoria2021violenciaenlinea}. La extensión a daño legal/económico y la fusión son decisiones locales. \\
\Sex & \texttt{sexual\_explicito}, \texttt{sexual\_cosificacion}, \texttt{sexual\_no\_consensual} & Descripción sexual gráfica sin fin informativo, reducción de una persona a objeto sexual o referencia a material íntimo sin consentimiento. No toda mención sexual es daño; se excluye el uso educativo, periodístico, documental, científico o artístico cuando el contexto lo sustenta. El sistema observa texto, no desnudez visual. & Tipología de daño, misoginia, violencia peruana en línea y política de plataforma \cite{banko2020taxonomy,zeinert2021misogyny,albornoz2018conocer,defensoria2021violenciaenlinea,youtube2026sexualpolicy}. \\
\texttt{SEGURO} & \texttt{seguro}, \texttt{seguro\_ironia\_marcada} & Categoría aprendida, mutuamente excluyente con daño. El segundo estado fino exige que la ironía tenga por blanco una institución, política o situación y no ataque a una persona o grupo protegido. Una lista de daños vacía no basta para asignarla. & Regla operacional del proyecto; no es una categoría experta ni legal. \\
Flags & \texttt{ironia\_ambigua}, \texttt{humor\_encubridor}, \texttt{contexto\_necesario} & Señalan sentido incierto, humor que acompaña daño o evidencia insuficiente en el chunk. Solo se conservan junto a un daño sustentado y activan revisión humana; no son clases de daño. & Abuso implícito, ironía, humor racializado y dependencia del contexto \cite{waseem2017abuse,elsherief2021implicit,ilic2018irony,salem2016amixer,bourgeade2024context,thakur2025quechua}. \\
\bottomrule
\end{tabularx}
\end{table*}

La carpeta de diseño contiene la entrada, el esquema de salida y 60 casos de verificación que cumplen el formato previsto. La guía 1.3 se conserva como parte del protocolo de anotación y la Tabla~\ref{tab:taxonomia} presenta por separado el sustento académico de las categorías.

En modelos clásicos y Transformers compactos, las finas y flags se reservaron para derivación o auditoría. En Qwen fueron objetivos auxiliares predichos desde texto, nunca variables \emph{gold} entregadas como entrada. Las reglas que limitan confianza a 0,65 ante ironía/contexto y activan revisión por debajo de 0,70 también pertenecen al protocolo local; no son umbrales tomados de la literatura.

\subsection{Preanotación y adjudicación}

El flujo completo combina preanotación, revisión dirigida y adjudicación. DeepSeek Flash realiza la preanotación amplia y DeepSeek Pro revisa de forma selectiva las propuestas con daño, baja confianza o flags, además de un control reproducible de seguros y negativos difíciles. Buscar rótulos posiblemente erróneos responde a un riesgo conocido \cite{brodley1999mislabeled}, pero la minería concreta de este proyecto es una auditoría heurística propia. Los casos todavía inciertos pasan a revisión final asistida.

La cola final presenta la propuesta de Pro durante la revisión y conserva la decisión resultante. La procedencia del corpus es 93\,636 chunks Flash (80\%), 20\,516 Pro o minería Pro (18\%) y 3\,092 revisiones humanas asistidas (2,6\%); los porcentajes se redondean de forma independiente. La precedencia es revisión humana final \(>\) Pro de minería \(>\) Pro \(>\) Flash.

El Apéndice~\ref{app:interfaces} muestra la vista final usada para contrastar el texto, su intervalo y la propuesta del LLM antes de aceptar, rechazar o modificar la decisión.

\subsection{Particiones y control contra fuga}

Para la comparación común se conservan todos los casos con daño y cuatro seguros elegidos por hash por cada caso positivo. La muestra resultante contiene 20\% de daño y, por ello, no representa la distribución natural de la plataforma \cite{huang2021balancing,tonneau2024naijahate}. Las particiones se forman por video con proporción 70/15/15, de modo que ningún video se comparte entre entrenamiento, validación y prueba \cite{sklearn2026groupshufflesplit}. Las puertas binarias incorporan negativos adicionales solo desde videos asignados a entrenamiento; las otras particiones no se reutilizan para ajustar modelos.
